\documentclass[journal]{IEEEtran}
\usepackage[T1]{fontenc}
\usepackage{graphicx}
\usepackage{url}
\usepackage{cite}
\title{A Practical Template for Manuscripts and Rebuttals}
\author{Template Author}
\begin{document}
\maketitle
\begin{abstract}
This paper is a working manual for an IEEE manuscript and rebuttal template.
It explains numbered submission stages, dependency-aware builds, and point-by-point reviewer responses.
The example produces submission-ready documents.
\end{abstract}
\begin{IEEEkeywords}
LaTeX, manuscript preparation, peer review, reproducible builds.
\end{IEEEkeywords}
\section{Introduction}
Keep each submission in a separate directory, use Git to record its history,
and build documents through Make.
The references in this manual are fictional examples, not research claims.
The introductory workflow follows the example guide by Alpha~\cite{guide}.
\section{Directory Structure}
Use \texttt{original} for the first submission, \texttt{revised1} and \texttt{reply1}
for the first review round, then \texttt{revised2} and \texttt{reply2}, and so on.
Every manuscript has \texttt{paper.tex} and \texttt{paper.bib}; section inputs and figures
remain local to that stage. Shared build helpers live at the repository root.
\begin{figure}[!b]
\centering
\setlength{\unitlength}{1mm}
\begin{picture}(80,22)
\put(0,12){\framebox(22,8){\texttt{original}}}
\put(23,16){\vector(1,0){10}}
\put(35,12){\framebox(25,8){\texttt{revised1}}}
\put(47,11){\vector(0,-1){3}}
\put(35,0){\framebox(25,7){\texttt{reply1}}}
\end{picture}

\caption{Each review round connects a revised manuscript to its response letters.}
\label{fig:workflow}
\end{figure}
\section{Installation and Builds}
Install GNU Make, Python 3, and TeX Live with IEEEtran, latexmk, latexpand, and latexdiff.
Use pdfunite for optional PDF bundles; Ghostscript is supported as a fallback.
Run \texttt{make check-tools} before the first build.
Alternatively, \texttt{make container} runs the same build in a pinned Linux container.
Run \texttt{make help} at the root or in a stage directory to discover its commands.
Root \texttt{make} builds all stages; \texttt{make -C original} builds the first manuscript.
All output is placed in ignored \texttt{build/} subdirectories.
\section{Comparing Manuscripts}
Every revision explicitly names its baseline in its Makefile.
For the first round this is \texttt{../original}; later rounds use the preceding revision.
Run \texttt{make -C revised1 diff} to build both manuscripts and their comparison.
The build flattens included source files and inserts each compiled bibliography
before marking changes. Blue underlining indicates additions and red strikeout indicates deletions.
This method follows the illustrative bibliography guide by Beta~\cite{bibliography}.
Generated comparison inputs always use \texttt{paper-expanded.tex}; the parent stage directory
identifies the version. Never rename these files or use them as manually maintained baselines.
A fictional automated-build checklist by Delta~\cite{automation} illustrates this approach.
Bibliography metadata corrections retain their citation keys. Removing a reference also removes
its citation from the manuscript; new references receive new keys.
The diff build regenerates its inputs automatically, including bibliography changes.

\section{Responding to Reviewers}
Edit \texttt{reply1/metadata.tex} to set the title, authors, and submission identifier.
Record each comment and answer in the corresponding reviewer file.
Individual and combined letters reuse the same answers.
Use \texttt{make -C reply1} to build the letters, or \texttt{make -C reply1 bundle}
to append the comparison to each letter.
\section{Adding a Review Round}
Run \texttt{make new-round} from the repository root.
This copies manuscript sources into the next revision, sets its baseline,
and creates fresh response placeholders. It skips generated files and refuses overwrites.
Root commands discover new stages automatically.
\section{Submission and Preservation}
Use \texttt{make -C revised1 zip} to create a standalone source archive.
Extract and compile the archive before uploading it.
Freeze submitted sources and tag them in Git, for example \texttt{submitted-revised1}.
Use \texttt{make clean} to remove auxiliaries while retaining PDFs and ZIPs;
\texttt{make distclean} removes all generated outputs.
\section{Conclusion}
Separate stages make the publication history easy to inspect.
Consistent commands and automatic baseline selection reduce the manual work needed to prepare each submission.
\bibliographystyle{IEEEtran}
\bibliography{paper}
\end{document}
